% Generated draft: requires mathematical and source review.
\documentclass[11pt]{article}
\usepackage[margin=25mm]{geometry}
\usepackage{fontspec}
\setmainfont{Latin Modern Roman}
\usepackage{amsmath,amssymb,mathtools,unicode-math}
\setmathfont{Latin Modern Math}
\usepackage{enumitem,needspace,xurl,hyperref}
\setlength{\parindent}{0pt}
\setlength{\parskip}{5pt}
\setlength{\emergencystretch}{3em}
\newcommand{\N}{\mathbb{N}}
\newcommand{\Z}{\mathbb{Z}}
\newcommand{\Q}{\mathbb{Q}}
\newcommand{\R}{\mathbb{R}}
\newcommand{\C}{\mathbb{C}}
\newcommand{\F}{\mathbb{F}}
\newcommand{\A}{\mathbb{A}}
\newcommand{\PP}{\mathbb{P}}
\newcommand{\E}{\mathbb{E}}
\newcommand{\eps}{\varepsilon}
\newcommand{\dd}{\,\mathrm{d}}
\begin{document}
\section*{UnsolvedMath ID: 20001958}
\textbf{Source problem code:} catalogue entry
\subsection{Definitions and mathematical statement}
G. Landweber asks whether every abstract moment map can be approximated by nondegenerate abstract moment maps, analogously to approximation of smooth functions by Morse functions. The source does not specify the action, topology, compactness or the nondegeneracy convention. The supplied negative answer uses the following standard torus convention.

Let a torus \(T\) with Lie algebra \(\mathfrak t\) act smoothly on \(M\). An abstract moment map is a smooth invariant map \(\Phi:M\to\mathfrak t^*\) such that for every subtorus \(H\subset T\), the projection \(\Phi|_{\mathfrak h}:M\to\mathfrak h^*\) is locally constant on \(M^H\). For \(\xi\in\mathfrak t\), write \(\Phi^\xi=\langle\Phi,\xi\rangle\) and \(\xi_M\) for the induced vector field. In the nondegenerate convention used here, every \(\Phi^\xi\) is Morse--Bott with
\[
 \operatorname{Crit}(\Phi^\xi)=\{x:\xi_M(x)=0\},
\]
and the Hessian kernel at a critical point is precisely the tangent space to that critical submanifold. The obstruction below already uses only the displayed critical-set requirement.

Suppose \(T=S^1\) acts locally freely on a nonempty compact manifold without boundary. Every invariant smooth real function is an abstract moment map: the only positive-dimensional circle subtorus has empty fixed set, so its fixed-set constancy requirement is vacuous. But for \(\xi\ne0\), the vector field \(\xi_M\) has no zero, whereas every smooth real function on this compact boundaryless manifold has a critical point. Thus no nondegenerate abstract moment map exists. In particular, for
\[
 M=T^2=\R^2/\Z^2,\qquad
 s\cdot(\theta,\phi)=(\theta+s,\phi),
\]
all invariant functions \(\Phi(\theta,\phi)=f(\phi)\) are abstract moment maps and none is nondegenerate. The nondegenerate locus is empty while the abstract-map space is nonempty, so density fails in every topology on this space. This does not contradict approximation by ordinary Morse functions, which are not required to satisfy the fixed action's nondegeneracy condition or invariance.

The aid also reports a necessary fixed-set condition. Let an \(r\)-torus act locally effectively on a compact connected boundaryless manifold, where local effectiveness means that the infinitesimal action has zero kernel. If a nondegenerate abstract moment map exists, then its values on connected components \(F\) of \(M^T\),
\[
 v_F=\Phi(F)\in\mathfrak t^*,
\]
affinely span \(\mathfrak t^*\). Consequently
\[
 \#\pi_0(M^T)\geq r+1.
\]
For clarity, fixed-component constancy comes from the abstract-map definition; generic infinitesimal components have critical set \(M^T\), so their extrema, and hence the convex support bounds for \(\Phi(M)\), occur at these values. A proper affine span would force some nonzero \(\Phi^\xi\) to be constant on all of \(M\), contradicting local effectiveness and the critical-set condition. This is a necessary condition, not a sufficient existence or density theorem. Noncompact manifolds, manifolds with boundary, restricted Hamiltonian settings or alternative definitions require separate formulations.
\subsection{Short English statement}
For a locally free circle action on a nonempty compact boundaryless manifold, invariant functions are abstract moment maps but none satisfies the required nondegeneracy condition. Thus the unrestricted density question fails. More generally, a locally effective rank-\(r\) action with such a map must have at least \(r+1\) fixed components.
\subsection{Sources}
\begin{itemize}
\item Permanent UnsolvedMath ID \(20001958\), exact public code catalogue entry; local registry \texttt{aim-geometry-notes.json}, \texttt{notes[295]}, canonical tag question.
\item AIM workshop \emph{Moment maps and surjectivity in various geometries}, Section 5 ``Topological aspects of moment map theory,'' Question 5.6, attributed to G. Landweber: \url{https://aimath.org/WWN/momentmaps/momentmaps.pdf}. No PDF was accessed.
\item Research attempt 1, ``Density fails by compactness for locally free circle actions,'' labelled counterexample, refuted at the run, proved and low novelty confidence; registry solved and SOLVED-IN-LITERATURE. No full authoritative literature citation, original literature entry or year is supplied for the definition or these results.
\end{itemize}
\end{document}
